Em avaliações de impacto estamos interessados em avaliar se um determinado programa ou intervenção foi efetivo. Utilizamos testes estatísticos para determinar se há diferenças nos resultados entre dois grupos (tratamento e controle). O uso desses testes se dá para nos ajudar a interpretar se as diferenças encontradas são apenas em virtude do acaso, erros de amostragem, ou se há evidências concretas de impacto. Em regra, os testes são feitos especificando uma hipótese nula de que as médias dos grupos não são diferentes umas das outras ou que a diferença média é zero.

A estatística t e o valor p nos ajudam a determinar se podemos rejeitar a hipótese nula com base no papel do "acaso" ou erro de amostragem na diferença observada entre as duas médias. Em testes estatísticos, podem ocorrer dois tipos principais de erros que influenciam nossos resultados: o falso positivo (\textbf{Erro Tipo I}) e o falso negativo (\textbf{Erro Tipo II}). Um falso positivo acontece quando rejeitamos a hipótese nula acreditando que há uma diferença ou um efeito significativo entre os grupos, quando, na verdade, essa hipótese nula é verdadeira. Em outras palavras, interpretamos que existe um impacto ou mudança, mas esse resultado é apenas um acaso do teste, sem efeito real. Esse erro pode nos levar a conclusões incorretas, interpretando variações aleatórias como efeitos importantes. Por outro lado, o falso negativo ocorre quando falhamos em rejeitar a hipótese nula, mesmo que exista um efeito real. Neste caso, o teste estatístico não detecta a mudança ou o impacto, levando-nos a concluir que não há efeito, quando na verdade ele existe. Esse erro nos impede de reconhecer uma diferença que deveria ser identificada.

Logo, para minimizar esses erros, nossa principal preocupação é determinar um tamanho de amostra grande o suficiente para restringir a prevalência de falsos positivos e falsos negativos a níveis mínimos aceitáveis na medição da diferença média "esperada" ou "assumida" entre duas populações. Esses parâmetros ajudam a tornar o teste mais confiável e preciso na detecção de impactos reais.

Como garantir que temos um tamanho de amostra suficiente para detectar um efeito, caso ele realmente exista? A resposta se dá através do Efeito Mínimo Detectável (\textbf{MDE}), que é a menor diferença entre as médias de dois grupos que conseguimos medir, considerando o tamanho da amostra, o poder do teste e a taxa de erro Tipo I definidos. Se a diferença observada for menor que o MDE, pode ser que não consigamos encontrar evidências estatísticas de um impacto, mesmo que ele realmente exista. Para calcular o tamanho amostral necessário, é importante considerar alguns fatores: o \textbf{nível de confiança}, que geralmente é fixado em 95\% para refletir o grau desejado de certeza; o poder do teste, que representa a probabilidade de detectar um efeito real e costuma ser definido em 80\% dependendo da dispersão/variância da variável de interesse do estudo.

Compreender o conceito de MDE é de extrema importância por diversas razões que afetam diretamente a condução e a eficácia de experimentos. Primeiramente, o MDE desempenha um papel crucial no planejamento do experimento, pois possibilita calcular o tamanho amostral necessário e definir a duração ideal para obter resultados confiáveis. Sem esse entendimento, corre-se o risco de planejar experimentos subdimensionados ou longos demais, comprometendo a eficiência e a validade das análises. Além disso, compreender o MDE é fundamental para uma alocação eficiente de recursos. Investir tempo, dinheiro e esforços em experimentos mal calibrados pode levar a desperdícios significativos ou à incapacidade de detectar efeitos reais. O MDE também oferece um contexto indispensável para a interpretação dos resultados obtidos nos experimentos. Ele estabelece uma referência clara sobre o tamanho mínimo do efeito que deve ser detectado, evitando interpretações equivocadas ou exageradas dos dados. Por fim, ter uma noção clara do MDE é essencial para gerenciar riscos de forma adequada, permitindo avaliar o impacto potencial das intervenções testadas e os possíveis riscos associados a elas. Dessa forma, o entendimento do MDE não é apenas uma ferramenta técnica, mas uma estratégia essencial para garantir o sucesso de experimentos bem planejados e bem executados.

O \textbf{Poder do teste} é um conceito intimamente relacionado à taxa de erro Tipo II. O poder é definido como 1 menos a taxa de erro Tipo II. Por exemplo, para uma taxa de erro Tipo II de 20\%, temos 80\% de poder. O poder é, portanto, a probabilidade de detectar um efeito quando esse efeito realmente existe.

Este guia tem como objetivo apresentar ao leitor o ferramental necessário para calcular o poder de experimentos. Na Seção \ref{poder}, primeiro derivamos o MDE para um caso padrão de Randomized Controlled Trials (RCT) e, a partir disso, discutimos como o cálculo do MDE muda para diferentes desenhos de experimentos nas seções subsequentes. A Seção \ref{late} irá tratar do cálculo de MDE para o caso de \textit{imperfect compliance}, a Seção \ref{cluster} discutirá o caso de Clustered RCT, a Seção \ref{strat} discutirá o caso de experimentos estratificados e a Seção \ref{multihip} irá tratar do caso em que queremos testar múltiplas hipóteses. Em todas essas seções, apresentaremos a teoria por trás da fórmula do MDE e traremos exemplos práticos utilizando o software R, calculando o MDE através de sua fórmula analítica e/ou por meio de simulações. 